% =====================================================================
% THESIS SKELETON: paragraph outline as LaTeX comments
% =====================================================================
%
% How the notes work (supervisor's advice, Nov 2026):
%   - Every paragraph has one block of comments before it.
%   - "P<n> KEY:" is the one idea of that paragraph. If a paragraph does not
%     serve its KEY, or its KEY does not serve the chapter's ROLE, cut it.
%   - "SUPPORT:" lists the facts, numbers and references the paragraph needs.
%   - "SRC:" points to where the numbers come from:
%       RO x.y  = "Results overview (2026-11).md", section x.y (OneDrive)
%       PL n    = PROGRESS_llms.md, section n (repo, branch main)
%   - "FIG:" / "TAB:" mark a planned figure or table.
%   - "TODO:" marks something not done or not checked yet.
%   - Keep the comments when the text is written. They help navigate.
%
% Length: official limit 35-70 pages for the whole thesis.
% Target: main text 40 pages, plus references and a short appendix.
% Page budget per chapter is in each chapter file.
%
% ---------------------------------------------------------------------
% MAIN STORY (every chapter must serve this)
% ---------------------------------------------------------------------
% Free LLMs can be prompted to write workplace email in a given role.
% (Q1) They show hierarchy, even more clearly than real Enron email does.
%      But the effect is fragile: it depends on the persona set, the
%      scenario and the model. It changes the wording and who does the
%      work, rarely the decision.
% (Q2) Their replies stay on topic, but they are easy to tell from real
%      email by style, and they are less diverse.
% (Q3) A strong LLM judge is consistent and does not favor its own
%      writing. But it needs the incoming message and has not been
%      checked against people.
% => LLM email simulation is useful for controlled "what if" comparisons
%    of hierarchy. It is not a stand-in for real organizational email,
%    and every result needs a real-email benchmark and a robustness check.
% =====================================================================

\documentclass[11pt,a4paper]{report}

\usepackage[utf8]{inputenc}
\usepackage[T1]{fontenc}
\usepackage{graphicx}
\usepackage{booktabs}
\usepackage{amsmath}
\usepackage[round]{natbib}
\usepackage{hyperref}

\title{LLMs for Simulating Organizational Communication and Decision-Making with Email}
\author{Euna Goo}
\date{TODO: submission date}

\begin{document}

\maketitle

\begin{abstract}
% KEY: the whole thesis in about 200 words. Write it last.
% SUPPORT (from the LMU SLDS template):
%   - the problem and the three questions
%   - the method: simulator, 4 free models, 3 designs, real-email benchmark
%   - the main results, with the most important numbers
%   - what the results mean
%   - no citations, no abbreviations, no figure references
\end{abstract}

\tableofcontents

% =====================================================================
% CHAPTER 1: INTRODUCTION                                   budget: 3 pages
% ROLE: Say why the question matters, state Q1-Q3, and preview the answer.
% Write this chapter last, when the results chapters are fixed.
% =====================================================================
\chapter{Introduction}
\label{ch:intro}

\section{Motivation}

% P1 KEY: Communication inside an organization shapes its decisions, and
%         hierarchy shapes that communication.
% SUPPORT:
%   - Exposé p.2: communication drives decision efficiency and alignment.
%   - Email is the main written channel in firms. Enron is the standard
%     public example.
%   - Power shows in workplace language (Gilbert 2012; Prabhakaran et al.).

% P2 KEY: LLMs are now used as stand-ins for people (as simulated agents)
%         and as evaluators of text (as judges).
% SUPPORT:
%   - Agents: Horton et al. ("Homo silicus"); Park et al. 2023 (Generative Agents).
%   - Judges: Zheng et al. 2023 (MT-Bench, LLM-as-a-judge).
%   - Promise: controlled experiments on role and hierarchy that are hard
%     to run with real employees.

% P3 KEY: Whether this works for organizational email is open: do simulated
%         replies show real hierarchy patterns, and can a judge score them?
% SUPPORT:
%   - Power-asymmetric conversations paper: LLM agents mirror some
%     socio-cognitive effects, not all.
%   - Moral Mazes (HR Simulator): LLM email tone differs from human email.
%   - Few studies compare simulated email to real email on the same measures.

\section{Research questions}

% P4 KEY: This thesis builds an email simulator with role personas and
%         asks three questions.
% SUPPORT:
%   - Q1: How does hierarchical role influence response content, sentiment
%     and decision attitude?
%   - Q2: Are LLM-generated responses diverse enough and similar to
%     empirical patterns?
%   - Q3: Can LLM-as-a-judge give consistent and calibrated scores for
%     communication quality?
%   - Wording from the Exposé (p.3). Keep it.

% P5 KEY: How each question is made measurable.
% SUPPORT:
%   - Q1: writing direction (up / to a peer / down) as the treatment.
%     Outcomes: orders, tone, politeness, decision field.
%   - Q2: generated vs. real replies to the same Enron emails, plus Self-BLEU.
%   - Q3: self-agreement, agreement between judges, self-preference, effect
%     of showing the incoming message, check against people.

% P6 KEY: The thesis uses only free models, which makes it a study of what
%         openly available LLMs can do.
% SUPPORT:
%   - No paid API (agreed with the supervisor).
%   - 4 models: Llama 3.2 3B (local), gpt-oss-20b, gpt-oss-120b, qwen3.8-27b
%     (free hosted tiers). Total cost $0.
%   - SRC: RO 0 (Models).

\section{Contributions}

% P7 KEY: What this thesis adds.
% SUPPORT:
%   - A reproducible simulator: 4 models, 3 designs, about 12,000 cached replies.
%   - A "real email first" rule: every hierarchy measure is tested on
%     real Enron email before it is used on generated replies.
%   - Q1: the hierarchy effect exists but is fragile (personas, scenario, model).
%   - Q2: the gap to real email is style, not topic.
%   - Q3: a corrected self-preference test (neutral reference writer) and
%     evidence that the judge needs the incoming message.

% P8 KEY: The answer in short (the main story from main.tex, in 3-4 sentences).
% SUPPORT:
%   - Copy the MAIN STORY block and shorten it once chapters 5-7 are written.

\section{Outline}

% P9 KEY: One sentence per chapter.
% SUPPORT:
%   - Ch. 2 background, Ch. 3 data and simulator, Ch. 4 measures,
%     Ch. 5-7 one chapter per question, Ch. 8 discussion and conclusion.

% =====================================================================
% CHAPTER 2: BACKGROUND AND RELATED WORK                    budget: 6 pages
% ROLE: Give the reader exactly the concepts that chapters 3-7 use:
%       (a) how hierarchy shows in text, (b) LLMs as simulated people,
%       (c) LLMs as judges, (d) how to compare generated and real text.
%       Anything not used later does not belong here.
% Literature notes on OneDrive: identifying_hierarchy.md,
%   authoritativeness-detection-text-research.md,
%   holistic-tone-authoritativeness-research.md,
%   llm-as-judge-communication-emotion-research.md,
%   llm-email-simulator-mirroring-research.md
% =====================================================================
\chapter{Background and Related Work}
\label{ch:background}

\section{Hierarchy and power in workplace language}
% budget: 2 pages

% P1 KEY: Formal rank and the power people show in text are not the same
%         thing. This thesis uses formal rank.
% SUPPORT:
%   - Bramsen et al. 2011 (Extracting Social Power Relationships).
%   - Prabhakaran et al. ("Who's (Really) the Boss?"): situational power
%     differs from the org chart.
%   - Nguyen et al. survey; Lam et al. 2018: text power and formal rank
%     often differ.
%   - Why it matters later: limitation in Ch. 8.

% P2 KEY: Research names several language markers of power. These are the
%         candidates for the Q1 measures.
% SUPPORT:
%   - Orders / controlling acts: Prabhakaran et al. (Overt Display of Power).
%   - Pronoun "you": Kacewicz et al. 2014.
%   - Hierarchy phrases: Gilbert 2012 (built on Enron).
%   - Hedges and commitment: Prabhakaran, Ganeshkumar & Rambow 2018.
%   - Politeness and gratitude: Danescu-Niculescu-Mizil et al. 2013.

% P3 KEY: Reading formal rank from one email is hard even for trained
%         models. So only a small real signal is expected.
% SUPPORT:
%   - Bramsen 2011 and Prabhakaran 2012: hierarchical power F1 34.8 on Enron.
%   - Lam et al. 2018 (Power Networks): a neural model for power relations.
%     TODO: take its reported scores from the paper.
%   - Ziems et al. (LLMs for CSS): zero-shot LLMs can classify power, with limits.

% P4 KEY: People also judge power from the whole message, not only from
%         single words. This motivates a holistic tone judge.
% SUPPORT:
%   - Körner et al. 2024: power is perceived from short texts, with consensus.
%   - Lens model: cue validity (does a cue go with the true trait?) and
%     cue utilization (does the reader use it?).
%   - Danescu-Niculescu-Mizil 2013: politeness as a human rating of the
%     whole message.

% P5 KEY: The Enron corpus is the standard data for this work, with a known
%         hierarchy.
% SUPPORT:
%   - About 500,000 emails from Enron employees (Exposé p.5).
%   - Agarwal et al. (Comprehensive Gold Standard for the Enron hierarchy).
%   - Used by most power studies above, so results can be compared.

\section{LLMs as simulated people}
% budget: 2 pages

% P6 KEY: LLMs given a persona can act as human-like agents in experiments,
%         with known caveats.
% SUPPORT:
%   - Horton et al. (Homo silicus): replicate economic experiments.
%   - Their caveats: memorization, "prompt hacking", causal inference is hard.
%   - "Using Large Language Models to Simulate Human Behavioural Experiments"
%     (Port of Mars game). TODO: check authors for the citation.

% P7 KEY: Agent simulations add memory and social settings.
% SUPPORT:
%   - Park et al. 2023 (Generative Agents): memory stream, retrieval, reflection.
%     Our persona memory follows this (Ch. 3).
%   - SOTOPIA: social goals and a holistic evaluation.
%   - Simulating Teams (VirT-Lab): team simulation and ground-truth checks.

% P8 KEY: The closest prior work studies power and workplace email
%         directly, but not against real email on the same measures.
% SUPPORT:
%   - "Do LLM Agents Mirror Socio-Cognitive Effects in Power-Asymmetric
%     Conversations": persona pairs like manager-employee.
%   - Moral Mazes in the Era of LLMs (HR Simulator): human vs. LLM email tone.
%   - Training tools (CommCoach, GLOSS) simulate workplace partners (Exposé p.7).

% P9 KEY: Simulators have known failure modes: echoing the input, agreeing
%         with the user, and repeating stock text.
% SUPPORT:
%   - Hao et al. (echo of prompt); Sharma et al. 2024 (sycophancy).
%   - Li et al. 2016 (dull, generic replies).
%   - Enron is likely in training data, so named personas could recall real
%     people. We use no names (Ch. 3).

\section{LLMs as judges}
% budget: 1.5 pages

% P10 KEY: An LLM judge scores text on a rubric. This is now common for
%          evaluating generated dialogue.
% SUPPORT:
%   - Zheng et al. 2023: pointwise scoring, MT-Bench.
%   - Gu et al. (survey): reliability has three axes, bias is one.
%   - Written score anchors (The Biased Oracle); several named dimensions
%     reported separately (ESC-Judge, HEART).
%   - Törnberg: LLMs can outperform expert coders at annotation.

% P11 KEY: Judges have known biases. The ones that matter here are
%          self-preference, length and rank.
% SUPPORT:
%   - Zheng et al. 2023: self-enhancement and verbosity bias. Advice: judge
%     from another model family.
%   - The Biased Oracle: rank bias, use more than one judge.

% P12 KEY: A judge must be checked before use: against itself, against other
%          judges, and against people.
% SUPPORT:
%   - Stability across repeated runs (voice-agent judge paper, arXiv:2608.24314).
%   - Human-coded samples as the standard check (ESC-Judge, HEART, MT-Bench).

\section{Comparing generated and real text}
% budget: 0.5 pages

% P13 KEY: Realism can be tested without a judge: by a classifier, by
%          embeddings, and by diversity.
% SUPPORT:
%   - Paired design: human and model replies to the same input (HEART).
%   - Style embeddings that ignore topic: StyleDistance (Patel et al. 2025).
%   - Self-BLEU (Zhu et al. 2018, Texygen), named in the Exposé.

\section{Research gap}
% budget: 0.5 pages

% P14 KEY: No study combines a role-based email simulator, a real-email
%          benchmark on the same measures, and a tested judge, using only
%          free models.
% SUPPORT:
%   - P8: prior work uses invented settings, no real benchmark.
%   - P11-P12: judge checks are rarely done with free, same-family judges.
%   - Leads to Q1-Q3.

% =====================================================================
% CHAPTER 3: DATA AND SIMULATION DESIGN                     budget: 6 pages
% ROLE: Describe the real data and the simulator well enough that a
%       reader could rebuild them. No results here.
% =====================================================================
\chapter{Data and Simulation Design}
\label{ch:design}

\section{Enron data and hierarchy labels}
% budget: 1.5 pages

% P1 KEY: How the raw corpus was turned into usable messages and threads.
% SUPPORT:
%   - Rebuild the message column, parse headers, find threads (Exposé p.8).
%   - TODO: fill in from PROGRESS.md / PROGRESS_llms.md (preprocessing
%     steps, duplicates removed, final counts).

% P2 KEY: Rank comes from job titles. Direction comes from comparing the
%         ranks of sender and recipient.
% SUPPORT:
%   - Job titles of 156 employees, 6 ranks; departments Trading and Legal.
%   - Direction: down / to a peer / up.
%   - Limitation (formal rank only): see Ch. 2 P1 and Ch. 8.
%   - SRC: RO 2.0.

% P3 KEY: Which real-email samples are used, and for what.
% SUPPORT:
%   - Q1 benchmark: 2,202 emails from 107 senders.
%   - Design A: 183 real threads (from 200 sampled).
%   - Direction judge (Layer 1): 1,001 real emails.
%   - Tone and politeness: 750 per source, 250 per direction.
%   - TAB: sample, size, used in which section.

\section{Models}
% budget: 1 page

% P4 KEY: Four models of different size and family, all free.
% SUPPORT:
%   - TAB: model, family, where it runs, role (writer / judge).
%   - Llama 3.2 3B local (Ollama); gpt-oss-20b, gpt-oss-120b, qwen3.8-27b
%     on NVIDIA and Groq free tiers.
%   - Why: no paid API; the laptop cannot run models above about 8B.
%   - SRC: RO 0 (Models).

% P5 KEY: Free tiers set practical limits, and a cache makes every result
%         reproducible.
% SUPPORT:
%   - Groq: about 1,000 requests and 200,000 tokens per day.
%   - Every reply is cached locally; a result never depends on the service.
%   - Total cost $0.
%   - Prompts with Enron text go to third-party servers (link to ethics, Ch. 8).

\section{Prompt structure}
% budget: 1.5 pages
% FIG: system diagram (update the Exposé diagram, p.9, to what was built).

% P6 KEY: Each persona is a role, built from real writing statistics, and
%         never a named person.
% SUPPORT:
%   - 10 personas: 5 ranks x 2 departments (Trading, Legal).
%   - Each summarizes how real employees at that rank wrote: length,
%     recipients, share of orders, hedges, deference, questions.
%   - Each built from at least 3 real people; rank 6 left out (too few).
%   - No names: the model may have memorized Enron (Ch. 2 P9).
%   - SRC: RO 1.1.

% P7 KEY: Each persona has a small memory, following Generative Agents.
% SUPPORT:
%   - 15 invented recent events (observations) and a few reflections.
%   - The most relevant ones for persona and direction go in the prompt.
%   - Park et al. 2023.

% P8 KEY: The reply comes as structured output. Its decision field is the
%         "decision attitude" of Q1.
% SUPPORT:
%   - 5 fields: subject, body, decision (accept / decline / defer /
%     escalate / none), confidence (low / medium / high), reason.

% P9 KEY: A hand reading found that small models hand requests back to the
%         sender, so the prompt says the persona must act.
% SUPPORT:
%   - About a quarter of 100 Llama replies mirrored the request.
%   - The added paragraph describes the behavior, not words, so it does not
%     target the mirroring measure.
%   - On Llama it changed little (borrowed words 0.579 -> 0.565, p=.61).
%   - Full prompt in the appendix.

\section{Three generation designs}
% budget: 1.5 pages
% TAB: the three designs side by side (input, personas, size, used for).

% P10 KEY: Design A answers real Enron emails, so each generated reply has a
%          real reply to compare with (Q2).
% SUPPORT:
%   - First message of a real thread; the persona matches the real replier's
%     rank and department; direction from the real ranks.
%   - 183 threads. Same paired logic as HEART.
%   - SRC: RO 1.2.

% P11 KEY: Design B is a full scenario grid that varies direction under
%          controlled conditions (Q1, Q3).
% SUPPORT:
%   - 6 task types x 3 directions x 2 stakes x 4 incoming tones = 144 scenarios.
%   - x 10 personas = 1,440 replies per draw. Llama: 3 draws (4,320).
%   - Incoming tone is a property of the message, not an instruction.
%   - Judge-swap set: 60 scenarios (2 task types x 3 directions, neutral tone).
%   - SRC: RO 1.3.

% P12 KEY: Design C uses 40 personas, so a rank effect cannot come from
%          one persona's prompt.
% SUPPORT:
%   - Problem: 10 personas = 10 clusters; clustered errors too small
%     (Cameron & Miller 2015); Llama's p moved from .248 to .018 between draws.
%   - 4 variants per role: tenure (2 or 12 years) x origin (internal / external).
%   - Each writes to a named receiver one rank up, same rank, one rank down.
%   - Pilot: 832 replies per model; qwen stopped (Groq limit).
%   - SRC: RO 1.4.

% P13 KEY: Generation worked for all models at no cost, but models write
%          shorter replies than people.
% SUPPORT:
%   - TAB: generation summary (RO 1.5): valid replies, mean length.
%   - Design A mean words: Llama 22.7, 20b 32.8, 120b 28.6, qwen 48.3, real 65.0.
%   - Length matters for several measures (Ch. 4, length matching).

\section{Changes from the original plan}
% budget: 0.5 pages

% P14 KEY: Some Exposé parts were dropped or changed, each for a reason.
% SUPPORT:
%   - Power score per message -> did not track rank; replaced by
%     sentence-level orders and judges.
%   - Scenario builder from threads -> Design A uses threads directly.
%   - Judge from another model family -> not possible on the free path;
%     replaced by the judge-swap design (Ch. 7).
%   - Human coders -> TODO: done or not? Decide before writing Ch. 7.

% =====================================================================
% CHAPTER 4: MEASURES AND STATISTICAL ANALYSIS              budget: 4 pages
% ROLE: Define every measure and model that chapters 5-7 report, once.
%       Results chapters then only point back here.
% =====================================================================
\chapter{Measures and Statistical Analysis}
\label{ch:measures}

\section{The real-email-first rule}
% budget: 0.5 pages

% P1 KEY: Every hierarchy measure is first tested on real email. A measure
%         that shows no hierarchy there is not used on generated replies.
% SUPPORT:
%   - Why: early simulator tests found no effect, with nothing to compare with.
%   - Dropped this way: "I", "we", apology, please, greeting, style matching.
%   - Same logic as Gupta et al. 2022; "cue validity" in Körner et al. 2024.
%   - SRC: RO 2.0, 2.3.

\section{Text measures}
% budget: 1 page

% P2 KEY: The main Q1 measure is the share of sentences that are orders.
% SUPPORT:
%   - Counted per sentence, because 57% of generated replies have one
%     sentence (a per-reply share is then 0 or 1).
%   - Imperative detection as in TV-AfD (precision 0.80).
%   - Validated power marker in Enron (Prabhakaran et al.).
%   - SRC: RO 2.1.

% P3 KEY: Secondary word measures, each with an expected direction from research.
% SUPPORT:
%   - "you" rate, hierarchy phrases (974 weighted phrases), gratitude, hedges.
%   - TAB (short): measure, expected direction, paper. SRC: RO 2.3.

% P4 KEY: Politeness comes from a model trained on human ratings.
% SUPPORT:
%   - xlm-roberta-large-tydip (Srinivasan & Choi 2022), trained on the
%     Stanford Politeness Corpus. Score: P(polite).
%   - Most emails score near 1, so the log-odds are analyzed.
%   - SRC: RO 2.6.

% P5 KEY: Two keyword markers show who does the work.
% SUPPORT:
%   - "Tells the other to act" (Send..., Pull..., Please...) and "I'll do it".
%   - Rough counts, checked by reading a few dozen replies. Not validated.
%   - SRC: RO 2.8.

\section{LLM judge instruments}
% budget: 1 page

% P6 KEY: Three judge tasks: guess the direction, rate the tone, score a rubric.
% SUPPORT:
%   - Direction (Layer 1): names removed, guess up / peer / down; score =
%     balanced accuracy (chance 33%). Layer 2: which cues go with "down".
%   - Tone: 1-7 on the whole email. Anchor 1 "defers to the reader",
%     anchor 7 "writes as the one in charge". No role shown.
%   - Rubric: 6 items, 1-5, with written anchors. Fidelity: role_consistency,
%     contextual_fit, corpus_plausibility. Performance: clarity,
%     politeness_appropriateness, conflict_management (Exposé p.8).
%   - Prompts in the appendix.

% P7 KEY: A judge is used only if it agrees with itself.
% SUPPORT:
%   - Weighted kappa over 3 passes on the same items; bar 0.4.
%   - Results of the gate are in Ch. 7 (Q3), the gate itself is used in Ch. 5.

\section{Realism and diversity measures}
% budget: 0.5 pages

% P8 KEY: Four measures compare generated and real replies to the same email.
% SUPPORT:
%   - Borrowed words: share of a reply's content words already in the incoming
%     email. Flag at >= 0.80. AUC 0.834 against 100 hand codes.
%   - Word classifier: TF-IDF + logistic regression, 5-fold CV, AUC.
%     Controls: length only; real reply cut to its partner's length.
%   - Embeddings: nomic-embed-text (topic + style) and StyleDistance (style only).
%     "Tracks its own email": closer to its own real partner than to another.
%   - Self-BLEU: vs. real replies (paired), and within one scenario.
%   - SRC: RO 3.1-3.5.

\section{Statistical models}
% budget: 1 page

% P9 KEY: Direction effects come from regressions that compare each
%         direction with writing to a peer.
% SUPPORT:
%   - Orders: logistic model, P(order) ~ direction + persona fixed effects,
%     errors clustered by persona.
%   - Real email: mixed model with a sender intercept and sender rank;
%     sender-clustered robust errors.
%   - Scenario and stakes analyses: OLS with persona fixed effects, HC3 errors.

% P10 KEY: Three checks guard against false positives.
% SUPPORT:
%   - Length matching: real emails cut to the model's mean reply length.
%   - Holm correction where many tests are run.
%   - Power analysis: e.g. a 0.6 tone-point gap needs about 55 replies per
%     direction at alpha=.05 (PL 26).
%   - Known weakness: 10 personas = 10 clusters (Ch. 3 P12).

% =====================================================================
% CHAPTER 5: Q1 - DOES HIERARCHY SHOW IN THE REPLIES?       budget: 8 pages
% ROLE: Answer Q1. Story of the chapter:
%   real email shows a small effect -> models show it too, even more ->
%   but it breaks with more personas and differs by scenario and model ->
%   it changes the wording and who does the work, rarely the decision.
% =====================================================================
\chapter{Q1: Does Hierarchy Show in the Replies?}
\label{ch:q1}

% P0 KEY: What this chapter tests and in which order.
% SUPPORT:
%   - Treatment: direction. Outcomes: orders (5.1), judge reads (5.2),
%     tone and politeness (5.3), scenarios and decisions (5.4).

\section{Orders in real email and in the simulator}
% budget: 2 pages

% P1 KEY: Real email shows a small hierarchy effect: more orders when
%         writing down.
% SUPPORT:
%   - Writing down vs. peer: +0.249 log-odds, p=.034. Writing up: none.
%   - Order share: 17.5% down, 14.2% peer, 15.0% up.
%   - Also follows research for "you" and hierarchy phrases. Gratitude goes
%     against it (more thanks when writing down).
%   - SRC: RO 2.1, 2.3.

% P2 KEY: With 10 personas, all four models show the same effect, often
%         larger, and give 2-3 times more orders overall.
% SUPPORT:
%   - Writing down: Llama +0.127, 20b +0.258, 120b +0.483, qwen +0.510 (all p<.001).
%   - gpt-oss-20b closest to real (p=.943 for the difference).
%   - FIG: q1_models_full_writing_down.png.
%   - SRC: RO 2.1.

% P3 KEY: Part of the gap to real email is reply length.
% SUPPORT:
%   - Real email cut to model length: no model differs from it (all p>.29).

% P4 KEY: With 40 personas, the writing-down effect disappears for every
%         model. It came from the few persona prompts, not from rank.
% SUPPORT:
%   - Llama -0.148 (p=.22), 20b +0.009 (p=.90), 120b -0.069 (p=.67).
%   - TAB: 10 vs. 40 personas per model.
%   - Caveat: 832 replies per model, less power than the 10-persona runs.
%   - SRC: RO 2.2.

% P5 KEY: The other word measures do not follow research in a stable way.
% SUPPORT:
%   - Models follow the expected direction in 3 of 16 cells (10 personas)
%     and 3 of 12 (40 personas).
%   - Only stable result: Llama hedges less writing down.
%   - Full table in the appendix. SRC: RO 2.3.

\section{Can a judge read the direction?}
% budget: 1 page

% P6 KEY: A judge cannot tell direction from one real email, but it can from
%         generated replies.
% SUPPORT:
%   - Only gpt-oss-120b passes the gate (kappa 0.542).
%   - Balanced accuracy: real 35.7% (chance 33%); Llama 35.6%; 20b 40.7%;
%     120b 43.5%; qwen 47.9%.
%   - FIG: role_inference_accuracy.png. SRC: RO 2.4.

% P7 KEY: The judge reads direction mainly from orders.
% SUPPORT:
%   - Orders are the main cue for every source (+0.43 to +1.03 log-odds per SD).
%   - "you" works in opposite directions in real and generated email.
%   - Caveat: 120b judges its own replies, but 20b and qwen also beat real email.

\section{Tone and politeness}
% budget: 2 pages

% P8 KEY: The larger models sound more in charge when writing down. Real
%         email and Llama do not.
% SUPPORT:
%   - Tone judge stable: kappa 0.877.
%   - Down vs. peer: real +0.24 (p=.17), Llama +0.07, 20b +0.41, 120b +0.45,
%     qwen +0.56 (last three p<.001). Only qwen also more deferent up (-0.42).
%   - Larger models sound more in charge overall (3.9-5.4 vs. 2.9-3.2 real).
%   - FIG: tone_judge_by_direction.png. SRC: RO 2.5.

% P9 KEY: A second judge confirms the large effects but not the small ones.
% SUPPORT:
%   - gpt-oss-20b as judge: weighted kappa 0.74 with 120b.
%   - Holds: 120b and qwen writing down, qwen writing up.
%   - Judge-dependent: 20b's own replies, Llama writing up.

% P10 KEY: Politeness changes with direction only for qwen. The gpt-oss
%          models give orders, but politely.
% SUPPORT:
%   - Median P(polite) 0.97-0.99 everywhere (ceiling).
%   - qwen: -0.45 down, +0.18 up (p<.001). Others and real email: no effect.
%   - Politeness goes against tone (Spearman -0.12 real, -0.25 to -0.45 models).
%   - SRC: RO 2.6.

% P11 KEY: For qwen, high stakes makes replies firmer in every direction,
%          not more deferent. The incoming tone matters little.
% SUPPORT:
%   - High stakes: tone +0.37 / +0.67 / +0.60 (down / peer / up), less polite.
%   - May follow the prompt wording ("real money and visible accountability").
%   - Assertive message: qwen matches a senior (+0.42) or peer, not a junior.
%     Weak result (no coefficient survives Holm).
%   - Tables in the appendix. SRC: RO 2.7.

\section{Scenarios and decisions}
% budget: 2 pages

% P12 KEY: A pooled effect can hide scenarios where direction matters a lot
%          and scenarios where it does not. So the effect is split by task type.
% SUPPORT:
%   - About 80 replies per scenario x direction cell. Holm inside each model.

% P13 KEY: qwen shows the gap in every scenario. It is about who does the
%          work: writing down it assigns tasks, writing up it takes them.
% SUPPORT:
%   - Tone gap from +0.29 (confirm_details) to +1.53 (schedule_coordination).
%   - resolve_disagreement: "tells other to act" 74% down vs. 29% up;
%     "I'll" 14% vs. 55%.
%   - One short example pair (from RO 2.8). Full table in the appendix.

% P14 KEY: The other models show the gap only in a few scenarios, each in
%          different ones. Llama shows it in none.
% SUPPORT:
%   - 120b: approve_or_decline (+1.32), report_problem (+0.68).
%   - 20b: report_problem (+0.49), resolve_disagreement (orders).
%   - So the scenario pattern belongs to the model, not to the task.
%   - FIG: scenario_direction_models.png. SRC: RO 2.8.

% P15 KEY: resolve_disagreement is the one shared case.
% SUPPORT:
%   - qwen and both gpt-oss models: more orders writing down, do the fix
%     themselves writing up. 20b "I'll" 45% -> 78%.

% P16 KEY: Only gpt-oss-120b lets rank change the decision itself.
% SUPPORT:
%   - approve_or_decline: down approve 21% / decline 40%; up approve 5% /
%     defer 66% (p<.001).
%   - Closest result to "deference toward a superior".
%   - Others: rank changes the wording, not the decision.
%   - Caveats: one draw; Llama's decision field repeats only about 60% of
%     the time; most small gaps are below the design's power.

\section{Summary}
% budget: 1 page

% P17 KEY: The answer to Q1 in three points.
% SUPPORT:
%   - Real email shows little hierarchy; models show more, mostly through orders.
%   - The effect is fragile: it breaks with 40 personas and depends on
%     scenario and model.
%   - Rank changes style and division of labor, rarely the decision.

% P18 KEY: Bridge to Q2: the replies show hierarchy, but do they look real?
% SUPPORT:
%   - Models give 2-3x more orders than people: a first sign of a style gap.

% =====================================================================
% CHAPTER 6: Q2 - DO GENERATED REPLIES LOOK LIKE REAL EMAIL? budget: 5 pages
% ROLE: Answer Q2. Story of the chapter:
%   models stay on topic -> but are easy to tell from real email, by style ->
%   judges still rate them as plausible -> and they are less diverse.
% =====================================================================
\chapter{Q2: Do Generated Replies Look Like Real Email?}
\label{ch:q2}

% P0 KEY: Q2 uses the paired Design A: real and generated replies to the
%         same 183 emails.
% SUPPORT:
%   - Exceptions: rubric (judge-swap set) and Self-BLEU within a scenario (Design B).

\section{Mirroring}
% budget: 0.5 pages

% P1 KEY: Copying the sender's words is mainly a weakness of the small model.
% SUPPORT:
%   - Borrowed words (same length): Llama 0.565 (19.1% flagged), 20b 0.373,
%     120b 0.435, qwen 0.298, real 0.294.
%   - All models borrow at least as much as real writers.
%   - FIG (optional): nvidia_four_models_mirroring.png. SRC: RO 3.1.

\section{Telling real from generated}
% budget: 2 pages

% P2 KEY: A simple word classifier tells real from generated replies almost
%         always.
% SUPPORT:
%   - AUC, same length: Llama 0.881, 20b 0.970, 120b 0.962, qwen 0.982.
%   - Length alone explains it only for Llama (0.837).
%   - SRC: RO 3.2.

% P3 KEY: The models are given away by stock phrases.
% SUPPORT:
%   - "I'll" in 59-61% of gpt-oss replies vs. 5-6% of real ones; "let"
%     (Let me know); "review"; "team".
%   - Real replies: sign-off names, "attached", "office".
%   - A larger model was not harder to spot.

% P4 KEY: In meaning space the replies stay tied to the email they answer.
% SUPPORT:
%   - nomic embeddings: AUC 0.82-0.94 (same length).
%   - "Tracks its own email" 0.81-0.86.
%   - SRC: RO 3.3.

% P5 KEY: Style alone separates them even better. The gap is how replies are
%         written, not what they say.
% SUPPORT:
%   - StyleDistance: AUC 0.97-0.99; "tracks its own email" falls to about 0.5,
%     as intended.
%   - Llama moves most (0.82 -> 0.97).
%   - FIG: embedding_style_vs_nomic_auc.png.

\section{Rubric scores from judges}
% budget: 1 page

% P6 KEY: Judges rate the larger models' replies as high quality and plausible.
% SUPPORT:
%   - Overall mean 4.4-4.6 of 5 under every judge.
%   - Llama depends on the judge (4.28 own judge; 3.48-3.60 gpt-oss judges),
%     mainly on role_consistency and contextual_fit.
%   - TAB: overall mean by writer x judge. Per-item table in the appendix.
%   - SRC: RO 3.4.

% P7 KEY: This disagrees with the classifier: "reads as plausible" is not the
%         same as "cannot be told apart from real".
% SUPPORT:
%   - Same replies: rubric ~4.5, classifier AUC ~0.97.
%   - The judge saw only the reply, not the role or message (Ch. 7).

\section{Diversity}
% budget: 1 page

% P8 KEY: All models are less diverse than real email.
% SUPPORT:
%   - Self-BLEU generated 0.51-0.57 vs. real 0.27-0.31 (same length), all p<.001.
%   - Fits P3: repeated stock phrases.
%   - FIG: self_bleu_vs_real.png. SRC: RO 3.5.

% P9 KEY: A different persona changes qwen's wording a lot and gpt-oss-120b's
%         very little.
% SUPPORT:
%   - Within one scenario: qwen 0.17, 20b 0.30, Llama 0.38, 120b 0.60.
%   - Model size does not decide diversity.

\section{Summary}
% budget: 0.5 pages

% P10 KEY: The answer to Q2, and the bridge to Q3.
% SUPPORT:
%   - On topic, but easy to spot by style, and less diverse.
%   - Judges call them plausible. So can the judge be trusted? -> Ch. 7.

% =====================================================================
% CHAPTER 7: Q3 - CAN AN LLM JUDGE BE TRUSTED?             budget: 4 pages
% ROLE: Answer Q3, and back up the judge results used in Ch. 5 and 6.
%   Story: consistent (on some questions) -> no self-preference once tested
%   correctly -> needs the incoming message -> not yet checked against people.
% =====================================================================
\chapter{Q3: Can an LLM Judge Be Trusted?}
\label{ch:q3}

% P0 KEY: Q1 and Q2 already use judges. This chapter checks them.
% SUPPORT:
%   - Every usable judge is also a writer (free path), so self-preference
%     must be tested.

\section{Consistency and agreement}
% budget: 1 page

% P1 KEY: Whether a judge agrees with itself depends on the question.
% SUPPORT:
%   - Direction: only gpt-oss-120b passes (0.542); 20b -0.016; small local
%     models fail.
%   - Tone: 120b 0.877, 20b 0.785, qwen 0.509 (all pass 0.4).
%   - TAB: judge x task kappa. SRC: RO 2.4, 2.5, 4.3.

% P2 KEY: Two judges agree on tone and find the same large effects.
% SUPPORT:
%   - Weighted kappa 0.74 between 120b and 20b; same score 39-45%, within
%     1 point 78-86%.
%   - FIG: tone_judge_two_judges.png.
%   - Caveat: both from the gpt-oss family.

\section{Self-preference}
% budget: 1.5 pages

% P3 KEY: The judge-swap design measures whether a judge scores its own
%         replies higher than expected.
% SUPPORT:
%   - All 4 models write replies to the same 60 scenarios; 3 models judge all.
%   - Self-preference = generator x judge interaction, against one reference
%     writer and one reference judge.
%   - Zheng et al. 2023 (self-enhancement bias).

% P4 KEY: The first reading found self-preference. It was wrong: the gpt-oss
%         judges are harsh on Llama, not kind to themselves.
% SUPPORT:
%   - Reference Llama: 120b +0.68, 20b +0.59 (p<.0001).
%   - Reference qwen (writes, does not judge): 120b -0.05 (p=.70),
%     20b -0.12 (p=.34).
%   - FIG: judge_swap_free_tier_self_preference_corrected.png. SRC: RO 4.1.

% P5 KEY: Lesson: report the interaction with a neutral reference writer and
%         with the cell means.

\section{Showing the judge the incoming message}
% budget: 1 page

% P6 KEY: Only gpt-oss-120b uses the incoming message well. With it, it
%         catches mirrored replies.
% SUPPORT:
%   - contextual_fit gap, sound vs. mirrored: 0.24 (p=.59) -> 0.65 (p=.001).
%   - 20b does not improve (0.37 -> 0.20).
%   - FIG: judge_context_contextual_fit_gpt_oss_120b.png. SRC: RO 4.2.

% P7 KEY: The message moves judges in opposite directions, and still shows
%         no self-preference.
% SUPPORT:
%   - Change in overall mean: Llama -0.25, 120b +0.21, 20b +0.30 (all p<.001).
%   - Scores with and without the message cannot be mixed.
%   - Self-preference with message vs. qwen: 120b +0.23 (p=.085), 20b +0.05.

\section{Calibration against people}
% budget: 0.5 pages

% P8 KEY: No person has checked any judge score yet.
% SUPPORT:
%   - Plan: 120 emails, 3 raters, about 1 hour each; would also check tone
%     and politeness.
%   - TODO: decide with the supervisor. If done, this becomes a result;
%     if not, a limitation in Ch. 8.

% P9 KEY: The answer to Q3.
% SUPPORT:
%   - gpt-oss-120b is consistent and not self-biased, but needs the incoming
%     message and has no human check.

% =====================================================================
% CHAPTER 8: DISCUSSION AND CONCLUSION                      budget: 4 pages
% ROLE: Put the three answers together, say what they mean, and be honest
%       about the limits. No new results here.
% =====================================================================
\chapter{Discussion and Conclusion}
\label{ch:discussion}

\section{The three answers together}
% budget: 1 page

% P1 KEY: Models show hierarchy more clearly than real people do, but the
%         effect is fragile.
% SUPPORT:
%   - Ch. 5: orders, judge reads, tone. Breaks with 40 personas.

% P2 KEY: The simulated hierarchy is a matter of style, not of decision.
% SUPPORT:
%   - Who does the work changes; the decision rarely does (only 120b).
%   - Fits Q2: the gap to real email is also style.

% P3 KEY: The judge is good enough to measure these effects, under conditions.
% SUPPORT:
%   - Ch. 7: gpt-oss-120b, with the message, large effects only.

\section{Implications}
% budget: 1 page

% P4 KEY: For LLM simulation research: the persona set is a design factor,
%         and a real-email benchmark is needed.
% SUPPORT:
%   - Homo silicus caveats (prompt hacking, memorization) seen here:
%     the stakes effect may follow the prompt wording.
%   - A result with 10 personas did not survive 40.
%   - Compare with the power-asymmetry paper and Moral Mazes.

% P5 KEY: For LLM judges: use a neutral reference writer, give the input,
%         and check against people.

\section{Limitations}
% budget: 1 page

% P6 KEY: Measurement limits.
% SUPPORT:
%   - No human coding (unless done, see Ch. 7).
%   - Rank from job titles, not perceived power.
%   - Keyword markers not validated; politeness model trained outside email.

% P7 KEY: Design limits.
% SUPPORT:
%   - 10 personas = 10 clusters; one draw for most models.
%   - Design C is a pilot (832 per model, no qwen).
%   - Stakes and tone effects only tested on qwen.

% P8 KEY: Model limits.
% SUPPORT:
%   - Free models only. Both reliable judges are from one family, and each
%     judges its own replies.

\section{Ethics}
% budget: 0.5 pages

% P9 KEY: Enron is real correspondence, so it was handled with care.
% SUPPORT:
%   - Analysis only; data not shared; no personas named after real people.
%   - Prompts with Enron text went to free third-party APIs.

\section{Conclusion and outlook}
% budget: 0.5 pages

% P10 KEY: The main story in one paragraph (see main.tex).

% P11 KEY: What comes next.
% SUPPORT:
%   - Human coding; the full 40-persona design; showing the judge the
%     writer's role while testing for rank bias.


\bibliographystyle{plainnat}
\bibliography{references}

\appendix
% =====================================================================
% APPENDIX                                          budget: about 8 pages
% ROLE: Material a reader needs to check or rebuild the work, but not to
%       follow the story. Keep it short: the page limit covers it too.
% =====================================================================

\chapter{Prompts}
\label{app:prompts}
% A1: full generation prompt, Design A (RO 1.2).
% A2: what changes in Designs B and C (RO 1.3, 1.4).
% A3: judge prompts: direction, tone, rubric with anchors.

\chapter{Personas and scenarios}
\label{app:personas}
% B1: the 10 personas (rank x department) and their writing statistics.
% B2: the 6 task types, stakes and incoming-tone texts.

\chapter{Additional results}
\label{app:results}
% C1: other word measures, 10 and 40 personas (RO 2.3).
% C2: stakes and incoming tone for qwen (RO 2.7).
% C3: qwen per-scenario table with keyword markers (RO 2.8).
% C4: rubric per item, all judges (RO 3.4).

\chapter{Reproducibility}
\label{app:repro}
% D1: code repository, how to rerun from the cache, manifests.
% D2: models, versions and service limits.


\end{document}
